\subsection*{The formalization of logic}

The general impression that we get from the (extremely incomplete) texts in our possession is that Greek philosophic thought of the fifth century B.C. was dominated by an increasingly conscious effort to extend to the whole field of human thought the procedures of articulate discourse which were put to use with such great success in rhetoric and contemporary mathematics --- in other words, to create Logic, in the most general sense of the word. The tone of philosophical writings underwent an abrupt change in this epoch. Whereas in the seventh and sixth centuries the philosophers asserted and prophesied (or, at best, sketched vague arguments, founded on equally vague analogies), from Parmenides and especially Zeno onward they reasoned and sought to extract general principles to serve as a basis for their dialectic. The first statement of the principle of the excluded middle is to be found in the work of Parmenides, and the proofs ``by contradiction'' of Zeno of Elea have remained famous. But Zeno flourished in the middle of the fifth century and, in spite of the uncertainty of documentation (*), it is very likely that at this period the mathematicians also used these principles in their own sphere.

As we have already said, it is not within the scope of this book to relate the innumerable difficulties which arose at each stage in the development of this Logic, and the resulting polemics which occupied the Eleatic school, and the Sophists, as well as Plato and Aristotle. Let us merely note the part played in this evolution by the assiduous culture of the art of oratory and the analysis of language which naturally followed from it --- developments which are attributed mainly to the Sophists of the fifth century. On the other hand, the influence of mathematics, if not always explicitly acknowledged, is no less clear, in particular in the writings of Plato and Aristotle. It could be said that Plato was almost obsessed by mathematics; although himself not an inventor in this domain, he kept himself informed of the discoveries of the contemporary mathematicians (many of whom were his friends or his pupils), and remained always most directly interested in the subject, to the extent of suggesting new directions for research. In his writings, it is mathematics which constantly serves as illustration or model (and on occasion, as with the Pythagoreans, feeds his penchant toward mysticism). Aristotle, as Plato's pupil, could hardly have failed to receive the minimum of mathematical education demanded of the pupils at the Academy, and a volume of selections from his writings which relate or allude to mathematics has been compiled {[}2 bis{]}. However, he seems never to have made much effort to keep in contact with the mathematical movements of his times, and in this domain he quotes only results which had long since become common knowledge. The later philosophers were, for the most part, even less well acquainted with mathematics. Many of them lacked the necessary technical knowledge, and their references to mathematics relate to stages long since passed in the evolution of the subject.

The culmination of this period, in so far as logic is concerned, is the monumental work of Aristotle {[}2{]}, whose great merit is to have succeeded for the first time in systematizing and codifying methods of reasoning which had hitherto remained vague or unformulated (*). For our purpose we need to keep in mind the general thesis of this work, namely that it is possible to reduce every correct argument to the systematic application of a small number of immutable rules which are independent of the particular nature of the objects under consideration (this independence is clearly brought out by the use of letters to denote concepts or propositions --- a usage which Aristotle very probably borrowed from the mathematicians). But Aristotle concentrated his attention almost exclusively on a particular type of relations and chains of reasoning, namely what he called ``syllogisms''; essentially, these are relations of the form A ⊂ B or A ∩ B ≠ ∅, in the language of set theory (†), and chains of such relations or their negations linked together by means of the scheme

\[
(\mathrm{A} \subset \mathrm{B} \text {   and   } \mathrm{B} \subset \mathrm{C}) \Longrightarrow (\mathrm{A} \subset \mathrm{C}).
\]

Aristotle was too well acquainted with the mathematics of his time not to have realized that schemes of this nature were not sufficient to account for all the logical operations used by mathematicians, nor a fortiori for other applications of logic ({[}2{]}, Analytica Priora, I, 35; {[}2 bis{]}, pp.~25-26) (*). At any rate, he undertook a close study of the various forms of ``syllogism'' (almost entirely devoted to the elucidation of the perpetual difficulties which arise from the ambiguity or obscurity of the terms involved in reasoning), which led him to formulate, among other things, rules for taking the negation of a proposition ({[}2{]}, Analytica Priora, I, 46). To Aristotle also belongs the merit of having distinguished with great clarity the role of ``universal'' propositions from that of ``particular'' propositions: a first step toward quantifiers (†). But it is only too well known how the influence of his writings --- often narrowly and unintelligently interpreted --- which remained considerable up to the nineteenth century, led philosophers to neglect the study of mathematics and impeded the progress of formal logic (**). Nevertheless, logic continued to make progress in antiquity, in the Megarian and Stoic schools, the rivals of the Peripatetics. Our knowledge of these doctrines is unfortunately all at second hand, often derived from their adversaries or mediocre commentators. It seems that the contribution of these logicians was the foundation of a ``propositional calculus'' in the modern sense of the word; instead of limiting themselves, like Aristotle, to propositions of the particular form A⊂B, they stated rules concerning entirely indeterminate propositions. Moreover, they analyzed so closely the logical relationships between these rules that they were able to deduce them all from five elementary rules, which were classed as ``unprovable'', by methods very similar to those we described in Chapter I, §§ 2 and 3 {[}5{]}. Unfortunately their influence was ephemeral, and their results lapsed into oblivion until rediscovered by the logicians of the nineteenth century. Aristotle remained the unchallenged master, in logic, until the seventeenth century. In particular, the scholastic philosophers were entirely under his influence, and although their contribution to formal logic is far from negligible, it contains no advance of the first order in comparison with that due to the philosophers of antiquity.

However, it does not appear that the works of Aristotle or his successors caused any notable repercussions in mathematics. The Greek mathematicians pursued their research along the paths opened by the Pythagoreans and their successors of the fourth century (Theodorus, Theaetetus, Eudoxus) without apparently concerning themselves with formal logic in the presentation of their results. This is hardly surprising when one compares the flexibility and precision of mathematical reasoning at this period with the extremely rudimentary state of Aristotelian logic. And, later, when logic passed beyond this stage, its evolution was guided by the new conquests of mathematics.

With the development of algebra, one could hardly fail to be struck by the similarity between the rules of formal logic and those of algebra, both being applicable to unspecified objects (propositions or numbers). And when, in the seventeenth century, algebraic notation took its definitive form in the hands of Vieta and Descartes, almost immediately there appeared various attempts to represent logical operations symbolically. But before Leibniz these attempts (as for example, that of Hérigone (1644) to write symbolically the proofs of elementary geometry, or that of Pell (1659) to do the same for arithmetic) were superficial and did not lead to any progress in the analysis of mathematical reasoning.

Leibniz was a philosopher who was also a mathematician of the first rank, and who was able to draw from his mathematical experience the germs of the ideas which were to release formal logic from the scholastic impasse (*). A universal genius if ever there was one, and an inexhaustible source of original and fertile ideas, Leibniz was all the more interested in logic because it lay at the heart of his grand projects of formalizing language and thought, at which he continued to work throughout his life. Having broken with scholastic logic already in his childhood, he was attracted by the idea (going back to Raymond Lulle) of a method of resolving all human concepts into primitive concepts, which thus constituted an ``alphabet of human thoughts'', and recombining them in a quasi-mechanical fashion so as to obtain all true propositions ({[}12 b{]}, vol.~VII, p.~185; cf.~{[}12 bis{]}, Chap. II). When still a very young man, he also conceived another, much more original idea, that of the utility of symbolic notation as a ``thread of Ariadne'' of thought (*). ``The true method'', he said ``should provide us with a filum Ariadnes, that is to say, a certain sensible material means which will lead the intellect, like the lines drawn in geometry and the forms of operations which are prescribed to novices in arithmetic. Without this, our minds will not be able to travel a long path without going astray'' ({[}12 b{]}, vol.~VII, p.~22; cf.~{[}12 bis{]}, p.~90). He had little knowledge of the mathematics of his time until the age of twenty-five, and first presented his projects in the form of a ``universal language''. But when he came into contact with algebra, he adopted it as the model for his ``universal characteristic''. By this phrase he meant a sort of symbolic language, capable of expressing all human thoughts unambiguously, of strengthening our powers of deduction, of avoiding errors by a purely mechanical effort of attention, and constructed in such a way that ``les chimeres, que celuy meme qui les avance n'entend pas, ne pourront pas estre ecrites en ces caracteres'' ({[}12 a{]}, vol.~I, p.~187). In countless passages in his writings where Leibniz alludes to this grandiose project and the progress which its achievement would bring (cf.~{[}12 bis{]}, Chapters IV and VI), one sees how clearly he conceived the notion of a formalized language, as a pure combination of signs in which only the order of writing is important (†), so that a machine would be capable of producing all the theorems (***) and all controversies could be resolved by a simple calculation ({[}12 bis{]}, vol.~VII, pp.~198-203). Although these hopes may appear extravagant, it is nonetheless true that this constant theme in Leibniz' thought underlies a good part of his mathematical output, beginning with his work on the symbolism of infinitesimal calculus. He himself was perfectly aware of this fact and explicitly associated his ideas on indicial notation and determinants ({[}12 a{]}, vol.~II, p.~240; cf.~{[}12 bis{]}, pp.~481-487) and his sketch of ``geometrical calculus'' (cf.~{[}12 bis{]}, Chapter IX) with his ``characteristic''. In his mind, the essential part of his scheme was to be symbolic logic or, in his words, a ``Calculus ratiocinator''; and if he did not succeed in creating this calculus, he made at least three attempts at it. In the first attempt he had the idea of associating with each ``primitive'' term a prime number, every term formed of several primitive terms being represented by the product of the corresponding prime numbers (*); he sought to translate into this system the usual rules of syllogism, but met considerable complications caused by negation (which he attempted, naturally enough, to represent by a change of sign) and soon abandoned this approach ({[}12 c{]}, pp.~42-96; cf.~{[}12 bis{]}, pp.~326-344). In later attempts, he sought to give Aristotelian logic a more algebraic form. Sometimes he retained the notation AB for the conjunction of two concepts, sometimes he used the notation A + B (†). He observed (in the multiplicative notation) the law of idempotence AA = A, noted that the proposition ``every A is B'' may be replaced by the equality A = AB, and that from this starting-point most of the rules of Aristotle may be obtained by purely algebraic calculation ({[}12 c{]}, pp.~229-237 and 356-399; cf.~{[}12 bis{]}, pp.~345-364). Also he had the idea of the empty concept (``non Ens'') and recognized, for example, the equivalence of the propositions ``every A is B'' and ``A. (not B) is not true'' (loc. cit.). Furthermore, he noted that his logical calculus applied not only to the logic of concepts but also to that of propositions ({[}12 c{]}, p.~377). Thus he came very close to ``Boolean calculus''. Unfortunately, he seems not to have succeeded completely in casting off the scholastic influence; not only did he limit the aim of his calculus almost entirely to the transcription of the rules of syllogism in his notation (**), but he went as far as sacrificing his most felicitous ideas to the desire of recovering the Aristotelian rules in their entirety, even those which are incompatible with the notion of the empty set (††).

The work of Leibniz remained, for the most part, unpublished until the beginning of the twentieth century, and had little direct influence. Throughout the eighteenth century and the beginning of the nineteenth, various authors (de Segner, Lambert, Ploucquet, Holland, De Castillon, Gergonne) gave outlines of attempts which were similar to those of Leibniz but which never really passed beyond the point at which he had halted. Their works made very little impression, and most of them were unaware of the results of their predecessors (*). Another who wrote under the same conditions was G. Boole, who can be considered as the true creator of modern symbolic logic {[}16{]}. His master idea was to take the ``extensional'' point of view systematically and thus to calculate directly with sets, denoting by \(xy\) the intersection of two sets, and by \(x + y\) their union when \(x\) and \(y\) are disjoint. He also introduced a ``universe'', denoted by 1 (the set of all objects), and the empty set, denoted by 0, and he wrote 1 --- \(x\) for the complement of \(x\) . As Leibniz had done, he interpreted the relation of inclusion by the relation \(xy = x\) (from which he easily deduced the justifications of the rules of the classical syllogism), and his notation for unions and complements gave his system a flexibility which its predecessors had lacked (†). Moreover, by associating with each proposition the set of ``cases'' in which it is satisfied, he interpreted the relation of implication as an inclusion, and his calculus of sets gave him in this way the rules of ``propositional calculus''.

In the second half of the nineteenth century, Boole's system became the basis of the work of an active school of logicians, who improved and completed it in various respects. Thus Jevons (1864) widened the sense of the operation of union \(x + y\) by extending it to the case where \(x\) and \(y\) are arbitrary. A. De Morgan in 1858, and C. S. Peirce in 1867, proved the duality relations

\[
\left(\mathbb {C} \mathrm{A}\right) \cap \left(\mathbb {C} \mathrm{B}\right) = \mathbb {C} (\mathrm{A} \cup \mathrm{B}), \quad \left(\mathbb {C} \mathrm{A}\right) \cup \left(\mathbb {C} \mathrm{B}\right) = \mathbb {C} (\mathrm{A} \cap \mathrm{B}) \quad (* *).
\]

In 1860 De Morgan also began the study of relations, and defined inversion and composition of binary relations (i.e., the operations which correspond to the operations \(\overline{\mathbf{G}}\) and \(\mathbf{G}_1 \circ \mathbf{G}_2\) on graphs) ( \(\ddagger\) ). All this work was systematically expounded and developed in the massive and prolix volumes of Schröder {[}23{]}. But it is rather curious to observe that the logicians we have just mentioned seemed scarcely interested in the application of their results to mathematics and that, on the contrary, the principal aim of Boole, and Schröder in particular, was apparently to develop ``Boolean'' algebra by imitating the methods and duplicating the problems of classical algebra (often in a rather artificial fashion). No doubt the reason for this attitude lies in the fact that Boolean calculus was still not suitable for transcribing the majority of mathematical arguments (*) and therefore was only a very incomplete answer to the grand dream of Leibniz. The construction of formalisms better adapted to mathematics --- the introduction of variables and quantifiers, due independently to Frege {[}24{]} and C. S. Peirce {[}22 b{]}, the capital step forward in this direction --- was the work of logicians and mathematicians who, in contrast to their predecessors, were primarily interested in applications to the foundations of mathematics.

Frege's project {[}24 b and c{]} was to base arithmetic on logic formalized by means of a symbolism to which he gave the name Begriffsschrift (concept-script); we shall return later to the method by which he defined the natural integers. His publications were characterized by extreme precision and minute care in the analysis of concepts, which led him to introduce many distinctions that have since shown themselves to be of great importance in modern logic. For example, he was the first to distinguish between the statement of a proposition and the assertion that the proposition is true; between the relation of membership and that of inclusion; between an object x and the set \(\{x\}\) whose only member is x, etc. This formalized logic, which contains not only ``variables'' in the mathematical sense of the word but also ``propositional variables'' which represent undetermined relations and are susceptible of quantification, was later (through the work of Russell and Whitehead) to provide the fundamental tool of metamathematics. Unfortunately the notation adopted by Frege was unsuggestive, of appalling typographical complexity, and far removed from current mathematical usage; and the resulting unreadability had the effect of considerably diminishing his influence on his contemporaries.

Peano's aim was both more ambitions and more down-to-earth. He set about publishing a Formulaire de mathematiques written entirely in formalized language, and containing not merely mathematical logic but all the results of the most important branches of mathematics. The speed with which he succeeded in completing this ambitious project, assisted by a constellation of enthusiastic collaborators (Vailati, Pieri, Padoa, Vacca,

Vivanti, Fano, Burali-Forti) bears witness to the excellence of the system he had adopted. Closely following current mathematical usage, and introducing many well-chosen abbreviating symbols, his language succeeded moreover in being fairly readable, thanks mainly to an ingenious system of replacing brackets by separating points {[}29 c{]}. Many of the notations due to Peano are now in common use by the majority of mathematicians: for example, \(\in\) , \(\supset\) (but, contrary to the present usage, in the sense of ``is contained'' or ``implies'' (*)), \(\cup\) , \(\cap\) , A---B (the set of differences \(a—b\) , where \(a\in A\) and \(b\in B\) ). Besides this, the Formulaire contains, for the first time, a close analysis of the general notion of a function, of direct ( \(\dagger\) ) and inverse images, and the remark that a sequence is just a function defined on N. But in Peano's hands quantification is subjected to irksome restrictions (in his system, in principle, only relations of the form A \(\Longrightarrow\) B, A \(\Longleftrightarrow\) B, or A = B can be quantified). Moreover, the almost fanatic zeal of some of his disciples laid him wide open to ridicule, and the often unjust criticism of H. Poincaré in particular was a serious blow to Peano's school and hindered the dissemination of his doctrines in the mathematical world.

With Frege and Peano we have the essential ingredients of the formalized languages in use today. The most famous of such languages is, without doubt, that created by Russell and Whitehead in their great work Principia Mathematica, which happily combines the precision of Frege with the convenience of Peano. Most of the formalized languages in current use differ from each other only in modifications of secondary importance whose object is to simplify their use. One of the most ingenious is the ``functional'' notation for relations (for example, \(\in xy\) in place of \(x \in y\) ), due to Lukasiewicz, which does away completely with the need for brackets. But the most interesting is certainly the introduction by Hilbert of the symbol \(\tau\) , which allows one to consider the quantifiers \(\exists\) and \(\forall\) as abbreviating symbols, to avoid the ``universal'' functional symbol \(\iota\) of Peano and Russell (which applies only to functional relations) and to dispense with the axiom of choice in the theory of sets ({[}31{]}, p.~183 ( \(^{**}\) )).

